\documentclass[12pt]{article}
\usepackage{amsmath, amssymb, bm}
\usepackage[margin=1in]{geometry}
\usepackage{physics}

\title{ME144 Project 2 \\ Foundational Models -- Dynamics \\ (Referenced to Zohdi v22)}
\author{Jack Nelson}
\date{}

\begin{document}
\maketitle

%%%%%%%%%%%%%%%%%%%%%%%%%%%%%%%%%%%%%%%%%%%%%%%%%%%%%%%%%%%%
\newpage
\section*{Problem 1: Single Particle Dynamics (Zohdi Sec. 2.1.2)}

\subsection*{1(a) Derive $a = v\,\dv{v}{x}$ using chain rule}

Start from the definitions (Zohdi Eq. (2.1)--(2.2)):
\[
v(t) = \dv{x}{t}, \qquad a(t) = \dv{v}{t}.
\]

Apply chain rule to $v=v(x(t))$:
\[
\dv{v}{t} = \dv{v}{x}\dv{x}{t}.
\]

Substitute $\dv{x}{t}=v$:
\[
a = \dv{v}{t} = \dv{v}{x}\,v.
\]

Therefore,
\[
\boxed{a = v\,\dv{v}{x}} \qquad \text{(matches Zohdi Eq. (2.6)/(2.8)).}
\]

\subsection*{1(b) With $x(0)=0$, $v(0)=0$: derive $x(t)$ under four acceleration forms}

\subsubsection*{(i) Constant acceleration $a=a_0$}

Start from definition:
\[
\dv{v}{t} = a_0.
\]
Separate and integrate from $t=0$ to $t$:
\[
\int_{v(0)}^{v(t)} dv = \int_0^{t} a_0\,dt.
\]
Using $v(0)=0$:
\[
v(t) - 0 = a_0 t \quad \Rightarrow \quad \boxed{v(t)=a_0 t}.
\]

Now use $v=\dv{x}{t}$:
\[
\dv{x}{t} = a_0 t.
\]
Integrate from $0$ to $t$:
\[
\int_{x(0)}^{x(t)} dx = \int_0^{t} a_0 s\,ds.
\]
Using $x(0)=0$:
\[
x(t) = a_0\left[\frac{s^2}{2}\right]_{0}^{t}
      = \frac12 a_0 t^2.
\]
So,
\[
\boxed{x(t)=\frac12 a_0 t^2} \qquad \text{(matches Zohdi Eq. (2.5) with $v_0=x_0=t_0=0$).}
\]

\newpage
\subsubsection*{(ii) Acceleration as a function of time $a=f(t)$}

Start:
\[
\dv{v}{t} = f(t).
\]
Integrate from $0$ to $t$:
\[
\int_{v(0)}^{v(t)} dv = \int_0^{t} f(\tau)\,d\tau.
\]
Using $v(0)=0$:
\[
\boxed{v(t) = \int_0^{t} f(\tau)\,d\tau}
\qquad \text{(Zohdi Eq. (2.10) with $v_0=t_0=0$).}
\]

Then,
\[
x(t)=\int_0^{t} v(s)\,ds
     =\int_0^{t}\left[\int_0^{s} f(\tau)\,d\tau\right]ds,
\]
so
\[
\boxed{x(t)=\int_0^{t}\left(\int_0^{s} f(\tau)\,d\tau\right)ds}
\qquad \text{(Zohdi Eq. (2.11) with $x_0=v_0=t_0=0$).}
\]

\newpage
\subsubsection*{(iii) Acceleration as a function of position $a=a(x)$}

Use the chain rule identity from part (a) (Zohdi Eq. (2.9)/(2.12)):
\[
a = v\dv{v}{x}.
\]
Rewrite:
\[
v\,dv = a(x)\,dx.
\]
Integrate from $(x,v)=(0,0)$ to $(x,v)=(x,v)$:
\[
\int_{0}^{v} v'\,dv' = \int_{0}^{x} a(\xi)\,d\xi.
\]
Compute the left integral:
\[
\left[\frac12 (v')^2\right]_{0}^{v} = \int_{0}^{x} a(\xi)\,d\xi
\quad\Rightarrow\quad
\frac12 v^2 = \int_{0}^{x} a(\xi)\,d\xi.
\]
Thus,
\[
v(x)=\sqrt{2\int_{0}^{x} a(\xi)\,d\xi}.
\]
Since $v=\dv{x}{t}$:
\[
\dv{x}{t}=\sqrt{2\int_{0}^{x} a(\xi)\,d\xi}
\quad\Rightarrow\quad
dt=\frac{dx}{\sqrt{2\int_{0}^{x} a(\xi)\,d\xi}}.
\]
Integrate from $0$ to $t$, and $0$ to $x(t)$:
\[
\boxed{
t=\int_{0}^{x(t)}\frac{d\eta}{\sqrt{2\int_{0}^{\eta} a(\xi)\,d\xi}}
}
\qquad \text{(matches Zohdi Eq. (2.12) logic, rearranged).}
\]

\newpage
\subsubsection*{(iv) Acceleration as a function of velocity $a=a(v)$}

Start:
\[
\dv{v}{t} = a(v).
\]
Separate:
\[
dt = \frac{dv}{a(v)}.
\]
Integrate from $0$ to $t$ and $0$ to $v(t)$:
\[
\boxed{
t=\int_{0}^{v(t)}\frac{d\mu}{a(\mu)}
}
\qquad \text{(Zohdi Sec. 2.1.2 item 4).}
\]

Also, use $v=\dv{x}{t}$ and substitute $dt$:
\[
dx = v\,dt = v\,\frac{dv}{a(v)}.
\]
Integrate:
\[
\boxed{
x(v)=\int_{0}^{v}\frac{\nu}{a(\nu)}\,d\nu
}
\qquad \text{(Zohdi Sec. 2.1.2 item 4).}
\]

%%%%%%%%%%%%%%%%%%%%%%%%%%%%%%%%%%%%%%%%%%%%%%%%%%%%%%%%%%%%
\newpage
\section*{Problem 2: Multidimensional Analysis (Zohdi Sec. 2.1.4 and Polar Eqns)}

\subsection*{2(a) Why do $\dot{\bm i}, \dot{\bm j}, \dot{\bm k}$ vanish in Cartesian coordinates?}

In a Cartesian frame, the basis vectors $\bm i,\bm j,\bm k$ are \emph{fixed in space} (they do not rotate with the motion),
so their time derivatives are zero:
\[
\boxed{\dot{\bm i}=\dot{\bm j}=\dot{\bm k}=0.}
\]

\subsection*{2(b) Normal--Tangential acceleration form}

In the moving tangent-normal basis,
\[
\boxed{\bm a = \dot v\,\bm e_t + \frac{v^2}{\rho}\,\bm e_n.}
\]
Here:
\begin{itemize}
\item $v=\abs{\bm v}$ is the speed,
\item $\dot v$ is the tangential acceleration magnitude,
\item $\rho$ is the local radius of curvature of the path,
\item $\bm e_t$ is the unit tangent direction of motion,
\item $\bm e_n$ is the inward unit normal direction (toward curvature center).
\end{itemize}

\newpage
\subsection*{2(c) Starting from Eq. (2.51), derive $\bm a$ in polar basis (use Eqns. (2.48), (2.50))}

From Zohdi Eq. (2.44), position is:
\[
\bm r = r\,\bm e_r.
\]
Velocity definition:
\[
\bm v = \dot{\bm r} = \frac{d}{dt}(r\bm e_r)
     = \dot r\,\bm e_r + r\,\dot{\bm e_r}.
\]
Use Zohdi Eq. (2.48):
\[
\dot{\bm e_r} = \dot\theta\,\bm e_\theta.
\]
Substitute:
\[
\bm v = \dot r\,\bm e_r + r\dot\theta\,\bm e_\theta,
\]
which is exactly Zohdi Eq. (2.51).

Now differentiate velocity to get acceleration:
\[
\bm a = \dot{\bm v}
= \frac{d}{dt}\Big(\dot r\,\bm e_r + r\dot\theta\,\bm e_\theta\Big).
\]
Apply product rule term-by-term:
\[
\bm a
= \frac{d}{dt}(\dot r\,\bm e_r) + \frac{d}{dt}(r\dot\theta\,\bm e_\theta)
= \ddot r\,\bm e_r + \dot r\,\dot{\bm e_r}
  + (\dot r\dot\theta + r\ddot\theta)\bm e_\theta + r\dot\theta\,\dot{\bm e_\theta}.
\]
Now substitute the polar basis derivatives:
\[
\dot{\bm e_r} = \dot\theta\,\bm e_\theta \quad \text{(Eq. 2.48)},\qquad
\dot{\bm e_\theta} = -\dot\theta\,\bm e_r \quad \text{(Eq. 2.50)}.
\]
So:
\[
\bm a
= \ddot r\,\bm e_r + \dot r(\dot\theta\,\bm e_\theta)
  + (\dot r\dot\theta + r\ddot\theta)\bm e_\theta
  + r\dot\theta(-\dot\theta\,\bm e_r).
\]
Combine like terms in $\bm e_r$ and $\bm e_\theta$:
\[
\bm a
= (\ddot r - r\dot\theta^2)\bm e_r
+ (2\dot r\dot\theta + r\ddot\theta)\bm e_\theta.
\]
Therefore,
\[
\boxed{\bm a = (\ddot r - r\dot\theta^2)\bm e_r + (r\ddot\theta + 2\dot r\dot\theta)\bm e_\theta}
\qquad \text{(matches Zohdi Eq. (2.52)).}
\]

%%%%%%%%%%%%%%%%%%%%%%%%%%%%%%%%%%%%%%%%%%%%%%%%%%%%%%%%%%%%
\newpage
\section*{Problem 3: Work--Energy Principle (Zohdi Sec. 2.3.5)}

\subsection*{3(a) Conservative force condition + demonstration with book example}

From Zohdi Sec. 2.3.5: if the work differential can be written as an exact differential,
\[
\bm F\cdot d\bm r = -dV,
\]
then
\[
U_{12} = \int_{\bm r_1}^{\bm r_2} \bm F\cdot d\bm r
       = -\int_{V_1}^{V_2} dV
       = -(V_2 - V_1).
\]
Also, since
\[
-dV = -\pdv{V}{x}dx - \pdv{V}{y}dy - \pdv{V}{z}dz,
\]
we identify
\[
F_x = -\pdv{V}{x},\quad F_y=-\pdv{V}{y},\quad F_z=-\pdv{V}{z}
\quad\Rightarrow\quad
\boxed{\bm F = -\nabla V.}
\]
Thus the defining conservative-force relationship is:
\[
\boxed{\bm F=-\nabla V.}
\]

\textbf{Demonstration using gravity (book example):}
For constant gravity, Zohdi lists
\[
\bm F = -mg\,\bm k \quad \Longleftrightarrow \quad V = mgz.
\]
Compute:
\[
\nabla V = \pdv{(mgz)}{x}\bm i + \pdv{(mgz)}{y}\bm j + \pdv{(mgz)}{z}\bm k
= 0\bm i + 0\bm j + mg\,\bm k,
\]
so
\[
-\nabla V = -mg\,\bm k = \bm F,
\]
which verifies the conservative relationship.

\subsection*{3(b) Meaning of $\int_1^2 \bm F_{\text{other}}\cdot d\bm r$}

\[
\boxed{
\int_{1}^{2} \bm F_{\text{other}}\cdot d\bm r
\text{ represents the work done by nonconservative (``other'') forces along the path from 1 to 2.}
}
\]

%%%%%%%%%%%%%%%%%%%%%%%%%%%%%%%%%%%%%%%%%%%%%%%%%%%%%%%%%%%%
\newpage
\section*{Problem 4: Conservation Conditions (Zohdi Sec. 2.4.1/2)}

\subsection*{4(a) Condition for linear momentum conservation in a given direction}

From Zohdi Sec. 2.4.1, linear momentum of a particle is
\[
\bm G = m\bm v,
\]
and the momentum form of Newton's 2nd law is (Eq. (2.89)):
\[
\sum \bm F_{\text{ext}} = \dot{\bm G}.
\]
If the net external force component in some direction ``dir'' is zero, then the corresponding
component of $\dot{\bm G}$ is zero:
\[
\sum F_{\text{ext,dir}} = 0 \;\Rightarrow\; \dot G_{\text{dir}}=0 \;\Rightarrow\; G_{\text{dir}}=\text{constant}.
\]
Therefore,
\[
\boxed{\text{Linear momentum is conserved in a direction if } \sum F_{\text{ext,dir}}=0.}
\]

\subsection*{4(b) Condition for angular momentum conservation about a point}

From Zohdi Sec. 2.4.2, angular momentum about a point $P$ is (Eq. (2.96))
\[
\bm H_P = \bm r_{P\to A}\times m\bm v,
\]
and the angular momentum balance has the form
\[
\sum \bm M_{P,\text{ext}} = \dot{\bm H}_P.
\]
If the net external moment about $P$ is zero, then
\[
\sum \bm M_{P,\text{ext}}=\bm 0 \;\Rightarrow\; \dot{\bm H}_P=\bm 0 \;\Rightarrow\; \bm H_P=\text{constant}.
\]
Therefore,
\[
\boxed{\text{Angular momentum about }P\text{ is conserved if }\sum \bm M_{P,\text{ext}}=\bm 0.}
\]

%%%%%%%%%%%%%%%%%%%%%%%%%%%%%%%%%%%%%%%%%%%%%%%%%%%%%%%%%%%%
\newpage
\section*{Problem 5: Kinetics of a System of Particles (Zohdi Sec. 2.5.1/2)}

\subsection*{5(a) Newton's 2nd law for a system of particles + meaning of terms}

From Zohdi Sec. 2.5.1: summing Newton's 2nd law over all particles and using cancellation of internal forces gives
\[
\sum \bm F_{\text{ext}} = \dv{}{t}\left(\sum m_i \bm v_i\right).
\]
Define the total linear momentum of the system:
\[
\bm P \equiv \sum m_i \bm v_i.
\]
Then
\[
\sum \bm F_{\text{ext}} = \dot{\bm P}.
\]
Using the center-of-mass (center-of-gravity) velocity identity $\bm P=M\bar{\bm v}_{cg}$ and differentiating (with $M$ constant),
\[
\dot{\bm P} = M\bar{\bm a}_{cg}.
\]
Thus the system form is (component forms are Eqs. (2.107)--(2.109), with the compact form consistent with Eq. (2.109)):
\[
\boxed{\sum \bm F_{\text{ext}} = \dot{\bm P} = M\bar{\bm a}_{cg}.}
\]
Interpretation:
\begin{itemize}
\item $\sum \bm F_{\text{ext}}$: sum of \emph{external} forces acting on the system,
\item $\bm P=\sum m_i\bm v_i$: total system linear momentum,
\item $M=\sum m_i$: total mass,
\item $\bar{\bm a}_{cg}$: acceleration of the center of mass (center of gravity under uniform gravity).
\end{itemize}

\subsection*{5(b) Is total KE equal to sum of particle KEs?}

Yes by definition:
\[
\boxed{T = \sum_{i=1}^{N}\frac12 m_i\,\bm v_i\cdot \bm v_i.}
\]

\newpage
\subsection*{5(c) Is total KE equal to KE of the center-of-mass? Explain}

Use Zohdi Eq. (2.111):
\[
\bm v_i = \bar{\bm v}_{cg} + \dot{\bm r}_{cg\to i}.
\]
Compute total kinetic energy:
\[
T = \sum \frac12 m_i \bm v_i\cdot \bm v_i
  = \sum \frac12 m_i (\bar{\bm v}_{cg} + \dot{\bm r}_{cg\to i})\cdot(\bar{\bm v}_{cg} + \dot{\bm r}_{cg\to i}).
\]
Expand the dot product:
\[
T
= \sum \frac12 m_i(\bar{\bm v}_{cg}\cdot\bar{\bm v}_{cg})
+ \sum \frac12 m_i(\dot{\bm r}_{cg\to i}\cdot\dot{\bm r}_{cg\to i})
+ \sum m_i(\bar{\bm v}_{cg}\cdot\dot{\bm r}_{cg\to i}).
\]
Rewrite each sum:
\[
T
= \frac12 (\bar{\bm v}_{cg}\cdot\bar{\bm v}_{cg})\sum m_i
+ \frac12 \sum m_i(\dot{\bm r}_{cg\to i}\cdot\dot{\bm r}_{cg\to i})
+ \bar{\bm v}_{cg}\cdot\sum m_i\dot{\bm r}_{cg\to i}.
\]
Use $\sum m_i = M$ and Zohdi’s fact that $\sum m_i \bm r_{cg\to i}=\bm 0$ implies $\sum m_i\dot{\bm r}_{cg\to i}=\bm 0$,
so the cross term vanishes. Therefore (Zohdi Eq. (2.113)):
\[
\boxed{
T = \frac12 M\bar{\bm v}_{cg}\cdot\bar{\bm v}_{cg}
  + \frac12 \sum m_i(\dot{\bm r}_{cg\to i}\cdot\dot{\bm r}_{cg\to i})
}
\]
The second term is generally nonzero, so in general
\[
\boxed{
T \neq \frac12 M\bar v_{cg}^2
}
\]
unless the particles have no relative motion about the center of mass.

%%%%%%%%%%%%%%%%%%%%%%%%%%%%%%%%%%%%%%%%%%%%%%%%%%%%%%%%%%%%
\newpage
\section*{Problem 6: Impulse and Momentum of a Particle System (Zohdi Sec. 2.5.3/4 style)}

\subsection*{6(a) Is COM momentum equal to sum of particle momenta?}

Start with the center-of-mass definition (Zohdi Eq. (2.103)):
\[
\bm r_{cg} = \frac{1}{M}\sum m_i \bm r_i.
\]
Differentiate:
\[
\bar{\bm v}_{cg} = \dot{\bm r}_{cg} = \frac{1}{M}\sum m_i \bm v_i
\quad\Rightarrow\quad
M\bar{\bm v}_{cg} = \sum m_i \bm v_i.
\]
But the total system momentum is (by definition, Sec. 2.5.3)
\[
\bm P \equiv \sum m_i \bm v_i.
\]
Therefore,
\[
\boxed{\bm P = \sum m_i\bm v_i = M\bar{\bm v}_{cg}.}
\]

\subsection*{6(b) Simple expression for conservation of angular momentum of a particle system}

If net external moment about point $P$ is zero:
\[
\sum \bm M_{P,\text{ext}} = \bm 0 \quad \Rightarrow \quad \dv{\bm H_P}{t}=\bm 0,
\]
so
\[
\boxed{\bm H_P = \text{constant}.}
\]

%%%%%%%%%%%%%%%%%%%%%%%%%%%%%%%%%%%%%%%%%%%%%%%%%%%%%%%%%%%%
\newpage
\section*{Problem 7: Rigid Body Dynamics (Zohdi Sec. 2.6.3 and 2.7.2/3)}

\subsection*{7(a) Using Fig. 2.26 and Eq. 2.152: why is normal velocity $v_n=0$ in pure rotation?}

In pure rotation about a fixed axis (Fig. 2.26), the velocity of a point $A$ is given by the rigid-body relation
\[
\bm v_A = \bm\omega \times \bm r_{O\to A}
\quad \text{(Eq. 2.152 per project hint).}
\]
Take the dot product with the radius vector $\bm r_{O\to A}$:
\[
\bm v_A\cdot \bm r_{O\to A}
= (\bm\omega\times \bm r_{O\to A})\cdot \bm r_{O\to A}.
\]
But for any vectors, $(\bm a\times \bm b)$ is perpendicular to $\bm b$, so
\[
(\bm\omega\times \bm r_{O\to A})\cdot \bm r_{O\to A}=0.
\]
Therefore the velocity has no component along the radial/normal direction (the direction of $\bm r_{O\to A}$):
\[
\boxed{v_n = 0.}
\]
(Equivalently: velocity is purely tangential in pure rotation.)

\subsection*{7(b) Summarize Eq. (2.190) in a sentence}

Zohdi Eq. (2.190) states:
\[
\sum \bm M_{\text{ext},p} = \dot{\bm H}_{cg} + \bm r_{p\to cg}\times M\bar{\bm a}_{cg}.
\]
In words: the net external moment about point $p$ equals the rate of change of angular momentum about the CG plus a “transport” term due to CG acceleration.

\newpage
\subsection*{7(c) Expand Eq. (2.196) and show planar motion reduces it to Eq. (2.197)}

Start from Zohdi Eq. (2.196):
\[
T = \frac12 M\bm v_{cg}\cdot \bm v_{cg}
  + \frac12 \{\bm\omega\}^T [I_{cg}] \{\bm\omega\}.
\]

\textbf{Planar motion assumption:}
rotation is only about the out-of-plane $z$ axis, so
\[
\{\bm\omega\} =
\begin{bmatrix}
0\\
0\\
\omega
\end{bmatrix}.
\]
Also in planar motion, Zohdi notes $[I_{cg}] \Rightarrow I_{cg,zz}$ (Eq. 2.186), i.e.
\[
[I_{cg}] =
\begin{bmatrix}
\ast & \ast & 0\\
\ast & \ast & 0\\
0 & 0 & I_{cg,zz}
\end{bmatrix}
\quad\Rightarrow\quad
\{\bm\omega\}^T [I_{cg}] \{\bm\omega\}
=
\begin{bmatrix}
0 & 0 & \omega
\end{bmatrix}
\begin{bmatrix}
\ast & \ast & 0\\
\ast & \ast & 0\\
0 & 0 & I_{cg,zz}
\end{bmatrix}
\begin{bmatrix}
0\\
0\\
\omega
\end{bmatrix}.
\]
Compute step-by-step:
\[
[I_{cg}]\{\bm\omega\}
=
\begin{bmatrix}
0\\
0\\
I_{cg,zz}\omega
\end{bmatrix},
\]
then
\[
\{\bm\omega\}^T([I_{cg}]\{\bm\omega\})
=
\begin{bmatrix}
0 & 0 & \omega
\end{bmatrix}
\begin{bmatrix}
0\\
0\\
I_{cg,zz}\omega
\end{bmatrix}
=
I_{cg,zz}\omega^2.
\]
Thus Eq. (2.196) becomes:
\[
T = \frac12 M v_{cg}^2 + \frac12 I_{cg,zz}\omega^2.
\]

Zohdi Eq. (2.197) also writes planar rigid-body KE in the decomposed/integral form:
\[
T = \frac12 M\bm v_{cg}\cdot\bm v_{cg} + \frac12 \int \dot{\bm d}\cdot\dot{\bm d}\,dm,
\]
and for planar rotation that second term evaluates to $\frac12 I_{cg}\omega^2$, giving:
\[
T = \frac12 M v_{cg}^2 + \frac12 I_{cg}\omega^2,
\]
consistent with Eq. (2.197) line following in the text.

\end{document}
